\documentclass[10pt,a4paper]{article}
\usepackage[margin=2cm]{geometry}
\usepackage[T1]{fontenc}
\usepackage{lmodern,booktabs,tabularx}
\usepackage[hidelinks]{hyperref}
\usepackage{microtype}
\setlength{\parindent}{0pt}
\setlength{\parskip}{5pt}
\title{VASP 6.6.1: UMT installation and validation record}
\author{MLIP-nucleation project}
\date{17 October 2026}
\begin{document}
\maketitle

\section{Installed software}
VASP 6.6.1 was compiled and validated as \texttt{u3129} on the UMT cluster.
The accepted CPU/MPI executables are \texttt{vasp\_std}, \texttt{vasp\_gam}
and \texttt{vasp\_ncl}. The installed root is
\path{/home/u3129/.local/opt/vasp/6.6.1-oneapi-2025}.
Executables are under \path{src/vasp.6.6.1/bin/} relative to that root.
Activate the installation by sourcing \path{env.sh} before an MPI launch
inside a suitable Slurm allocation. No global shell initialization file was modified.

The toolchain comprises Intel ifx 2025.3.2, Intel MPI 2021.17 and MKL 2025.3.
The build uses the supplied oneAPI template with project overrides:
\texttt{mpiifx}, sequential MKL and the CPU target \texttt{-xSSE4.2}.
GPU support, OpenMP, HDF5 and optional external interfaces were not enabled.
The supported scope is the CPU/MPI build tested on \texttt{apollo31};
performance tuning and other node/toolchain combinations are not claimed validated.

For \texttt{vdw\_nl.o}, \texttt{-fp-model precise} was applied and confirmed
in the compilation log. This follows ASU Research Computing's reported workaround
for the VASP 6.6.x/ifx nonlocal-vdW kernel optimization issue, accessed on
17 October 2026:\\
\url{https://docs.rc.asu.edu/vasp/}.

\section{Source and reproducibility}
The user supplied \texttt{vasp.6.6.1.tgz}, size 85,036,103 bytes.
Its locally computed SHA-256, verified again after transfer, is:\\
\path{261f9c1fd730468106191afeb185821d67cb7a58c982f6e063824dd5432106b0}.
This records the supplied archive and transport integrity, not an independent
comparison with the VASP Portal's release checksum.

Licensed sources, binaries and bundled test potentials remain in private
user directories on the cluster. The Git repository stores only project scripts,
settings, hashes and validation output. The accepted input manifest is
\path{runs/vasp661-oneapi-20261017-05/manifest.json}; input file hashes bind the
batch script, environment and compiler overrides to the actual submission.
Individual executable SHA-256 values and deployment checks are recorded in
\path{runs/vasp661-oneapi-20261017-05/result.json}.
Copied executables matched the tested executables byte for byte.

\section{Slurm build and tests}
Job \texttt{20950123} ran on \texttt{apollo31}, partition \texttt{apollo},
account \texttt{imet-2026-2}, QoS \texttt{normal}.
The resource cap was one node, eight CPUs, 16 GiB and one hour.
The successful job completed in 9 min 27 s with scheduler state
\texttt{COMPLETED} and exit code \texttt{0:0}.
Each variant was built with \texttt{make DEPS=1 -j8}; runtime library checks
found no unresolved shared dependencies.

Four explicitly selected upstream cases ran with four MPI ranks and one
MKL thread per rank. No additional scientific production calculations were run.
\begin{tabularx}{\linewidth}{@{}lXl@{}}
\toprule
Case & Coverage & Result \\
\midrule
\texttt{bulk\_AlPAW} & Standard metallic PAW calculation & Passed \\
\texttt{bulk\_BN\_vdW-DF3-opt1} & Nonlocal vdW functionality & Passed \\
\texttt{Fe\_r2SCAN\_ncl\_soc} & Noncollinear r2SCAN with SOC & Passed \\
\texttt{H2Ostretch\_efor} & Gamma-only calculation and force/trajectory checks & Passed \\
\bottomrule
\end{tabularx}

The testsuite reported \texttt{SUCCESS: ALL SELECTED TESTS PASSED}; test execution
took 1 min 3 s. Reference comparisons include energies, forces, stress and
other quantities as specified by the upstream recipes.
The full suite was not run. Successful installation tests do not replace
system-specific DFT convergence and accuracy studies.
The upstream test protocol is described at
\url{https://vasp.at/wiki/Validation_tests}.

The retained validation output is
\path{runs/vasp661-oneapi-20261017-05/validation-summary.txt};
runtime dependency records are in \path{runtime-libraries.txt} in that directory.
The bundled testsuite prints an unavailable Git SHA because the source was
distributed as an archive; the archive hash above identifies the source instead.

\section{Cluster-specific execution requirements}
Main home was readable but read-only on compute node \texttt{apollo20}.
Initial launches failed before compilation when attempting to create shared
stdout or runtime logs there. A Slurm diagnostic confirmed the read-only error;
a separate diagnostic on \texttt{apollo31} verified write access to
\path{/home6/u3129}. The build and tests therefore ran under
\path{/home6/u3129/.local/opt/vasp/6.6.1-oneapi-2025}.
Accepted executables were copied to main home from the login node after validation.
Calculation working directories and Slurm logs must use compute-writable storage.

Compute-node Lmod did not expose the login module
\texttt{intelmpi/intel64\_2025}. The installation environment instead sources
the inspected shared compiler, MPI and MKL initialization scripts directly,
with pinned component versions. Existing user MLIP environments were not modified.

The supplied test POTCAR files were sufficient for the selected validation
cases. A full licensed PAW potential dataset and its provenance must still be
configured before production DFT labelling. No general production potential
library is inferred from these tests. Storage backup and retention policies
remain to be confirmed for the research workflow.

\end{document}
